% C.3bm. Complete root proof and both independent whole-proof audits passed.
\paragraph{C.3bm. The universal exceptional-surface reduction in positive
characteristic, dimension three (Proved reduction).}
Let $k$ be any algebraically closed field of characteristic $p>0$.
Let $R$ be the strict henselization at a closed $k$-point of a normal
local ring essentially of finite type over $k$. Suppose that
$\dim R=3$ and that
\[
             X=\operatorname{Spec}R\setminus\{\mathfrak m_R\}
\]
is regular. There exists a regular projective modification
$f:Z\to\operatorname{Spec}R$, an isomorphism over $X$, whose full
reduced closed fibre $E=\bigcup_{i=1}^c E_i$ is an SNC divisor and
which satisfies the hypotheses of EL and EP. In particular, with
\[
 S(E)=F_\gamma^2K_0(E)_{\mathbf Q},\qquad
 J_K(E)=\left\langle\beta_t-\beta_o:
 \begin{array}{l}
 P/k\text{ smooth connected projective curve},\\
 o,t\in P(k),\quad\beta\in F_\gamma^2K_0(E\times P)_{\mathbf Q}
 \end{array}\right\rangle,
\]
and with $\theta:\mathrm{CH}_1(E)_{\mathbf Q}\to S(E)$ the
exceptional pushforward of EP, there is an isomorphism
\begin{equation}
 B^2(X)_{\mathbf Q}\simeq
 \frac{S(E)}{J_K(E)+\theta\mathrm{CH}_1(E)_{\mathbf Q}}.
                                                    \tag{C.3bm.1}
\end{equation}
Thus the exceptional-family classes defining the relation subgroup
are unrestricted actual perfect classes on $E\times P$: no
extension of that class to $Z\times P$, and no vanishing of its
fixed-family negative-$K$ boundary, is required.

There is also an explicit degree-zero form of the remaining gate.
Define
\[
 \deg_E:S(E)\longrightarrow\mathbf Q^c,
 \qquad a\longmapsto
       \bigl(\deg\operatorname{ch}_2(a|_{E_i})\bigr)_i,
 \qquad S_0(E)=\ker(\deg_E).
\]
Then $J_K(E)\subset S_0(E)$, and
\begin{equation}
 B^2(X)_{\mathbf Q}\simeq
 \frac{S_0(E)}{J_K(E)+
        \bigl(S_0(E)\cap\theta\mathrm{CH}_1(E)_{\mathbf Q}\bigr)}.
                                                    \tag{C.3bm.2}
\end{equation}
The theorem holds for every prime, including $2$ and $3$.
It removes the ambient endpoint-lifting condition from this entire
algebraizable threefold branch. It does not assert that the quotient
in (C.3bm.2) vanishes or is finite-dimensional, and makes no claim
in characteristic zero.

\begin{proof}
\emph{1. The required geometry over an arbitrary algebraically closed field.}
We spell out why the construction used in C.3bd is independent of
its special coefficient field. If $R$ is regular, blow up its closed
point. Then $Z$ is regular, $E=\mathbf P^2_k$, and $D=E$ has
$\mathcal O_E(-E)=\mathcal O_{\mathbf P^2}(1)$.

Suppose that $R$ is singular. Apply
\citet[Theorem 1.1, author pp.~3--4]{frC3bdCossartPiltant2019}
to the affine excellent threefold $\operatorname{Spec}R$.
It provides a projective regular modification
$\pi_1:Z_1\to\operatorname{Spec}R$ which is an isomorphism over
$X$. Its projectivity follows from the theorem's specified-affine-cover
clause; it is not inferred merely from properness. Choose a
$\pi_1$-ample line bundle $H_1$, set
\[
 \mathcal L_1=(\pi_1)_*H_1,
 \qquad\mathcal L_2=\mathcal Hom(\mathcal L_1,\mathcal O_R),
 \qquad Z_2=\operatorname{Proj}_R
                 \bigoplus_{n\ge0}\mathcal L_2^{\otimes n}.
\]
Both sheaves are generically of rank one and are invertible on
$X$; hence $Z_2\to\operatorname{Spec}R$ is the identity there.
On the regular $Z_1$, the torsion-free quotient $M$ of
$\pi_1^*\mathcal L_2$ embeds into its invertible reflexive hull.
The coherent ideal
\[
                 I_1=M\otimes(M^{**})^{-1}\subset\mathcal O_{Z_1}
\]
is the unit ideal on $X$ and at every codimension-one point.
\citet[Proposition 4.4 and its proof, author p.~141]
{frC3bdCossartPiltant2019} principalizes $I_1$ by a finite sequence
of blow-ups in regular centres confined to the nonprincipal loci.
The resulting regular $Z_0$ remains projective and unchanged over
$X$. The now invertible quotient of the pullback of $\mathcal L_2$
gives a morphism $Z_0\to Z_2$.

Apply \citet[Lemma 8 and its proof, author pp.~2--3]
{frC3bdKollarWitaszek2026} to the actual common model $Z_0$ over
$Z_1$ and $Z_2$. It supplies an ample exceptional Cartier divisor
$F_0$ for $g_0:Z_0\to\operatorname{Spec}R$.
Every exceptional centre is over the closed point. The full reduced
closed fibre has dimension at most two. Therefore
\citet[Corollary 0.4 and its proof, author p.~5]
{frC3bdCossartJannsenSaito2013} supplies a regular projective
$h:Z\to Z_0$, unchanged outside this fibre, for which its entire
reduced inverse image is an SNC divisor $E$.
Since $Z_0$ is regular and thus factorial locally,
Koll\'ar--Witaszek's Lemma~6 gives an $h$-ample exceptional
Cartier divisor $G$. After taking a positive multiple if necessary,
\[
                         F=h^*(mF_0)+G
\]
is $f=g_0h$-ample for sufficiently large $m$, by the relative
ampleness composition used also in their Lemma~8. Its support
is entirely over the closed point, and $f$ remains an isomorphism
over $X$.

\emph{2. Positivity, full support, and the nonreduced divisor.}
For sufficiently large $n$, $\mathcal O_Z(nF)$ is globally generated
because the base is affine and it is relatively ample
\citep[Proposition 28.27.13(7)]{frC3bdStacksGeneration}.
Use its rational-function trivialization over $X$. Normal Hartogs
\citep[Remark 31.17.2]{frC3bdStacksHartogs} identifies
\[
                 I=f_*\mathcal O_Z(nF)\subset R
\]
with a coherent ideal that is the unit ideal off the closed point.
The evaluation surjection gives, inside the common function field,
\[
                 I\mathcal O_Z=\mathcal O_Z(nF)\subset\mathcal O_Z.
\]
Thus $D=-nF$ is effective. The ideal $I$ is proper: otherwise the
relatively ample line bundle would be trivial, every projective
fibre would be zero-dimensional, and the resulting finite birational
map to the normal local scheme would be an isomorphism. This would
make its singular closed point regular. Hence $I$ is
$\mathfrak m_R$-primary. Its pullback vanishes on the full closed
fibre, so
\[
                           D=\sum_i a_iE_i,\qquad a_i>0.
\]
The line bundle $\mathcal O_E(-D)$ is ample. Since $E\hookrightarrow D$
is a nilpotent thickening, ampleness also holds on $D$ itself by
\citet[Lemma 30.17.5, Tag 09MS]{rootStacks2026NilpotentAmple}. In particular
$\mathcal O_D(-D)$ has precisely the ampleness required by EL.
All of these arguments use only excellence, normality, the
three-dimensional resolution theorems and projectivity; none uses
that $k$ is algebraic over a finite field.

\emph{3. Apply the proved endpoint comparison.}
The constructed data satisfy all hypotheses of EL and EP.
Finite projective data over the strict localization descend to
pointed etale finite-type models as in EP, so its Chern-character
comparison applies. In particular
\[
 \rho:\mathrm{CH}^2(Z)_{\mathbf Q}\xrightarrow{\sim}S(E),
 \qquad
 \rho(\mathrm{CH}^2(Z)_{\mathrm{alg},\mathbf Q})=J_K(E).
\]
The second equality is the substantive endpoint theorem: every
perfect gamma-two family on the whole exceptional product has
an actual family on the fixed $Z\times P$ with the same two
endpoints, even if the original whole class does not extend.
Ordinary Chow localization, with closure of actual parameter
families as proved in EP, gives (C.3bm.1).

\emph{4. Remove all component-degree coordinates.}
For every $\beta$ used in $J_K(E)$ and every component $E_i$,
proper pushforward of $\operatorname{ch}_2(\beta|_{E_i\times P})$
to the connected curve $P$ has one constant degree. Refined
pullback at $o$ and $t$ preserves this degree. Therefore
$J_K(E)\subset S_0(E)$.

Choose an ample line bundle $H$ on $E$ and put
$C_j=c_1(H|_{E_j})\cap[E_j]\in\mathrm{CH}_1(E_j)$.
The degree vector of $\theta(C_j)$ is column $j$ of the symmetric
matrix
\[
       M_{ij}=\deg\bigl(c_1(\mathcal O_Z(E_i)|_{E_j})
                                  c_1(H|_{E_j})\bigr).
\]
For $i\ne j$ this is a nonnegative degree on $E_i\cap E_j$.
With $a=(a_i)_i$ from Step~2, ampleness gives $(Ma)_i<0$.
Put $A=\operatorname{diag}(a_i)$ and $B=AMA$. Its off-diagonal
entries are nonnegative and its row sums are strictly negative.
For every nonzero real vector $x$,
\[
 x^{\mathsf T}Bx
   =\sum_i\Bigl(\sum_jB_{ij}\Bigr)x_i^2
       -\sum_{i<j}B_{ij}(x_i-x_j)^2<0.
\]
Thus $M$ is invertible over $\mathbf Q$. In particular exceptional
cycles span all degree coordinates. Explicitly, for $s\in S(E)$,
\[
 s-\sum_j\bigl(M^{-1}\deg_E(s)\bigr)_j\,\theta(C_j)\in S_0(E).
\]
Every coset in (C.3bm.1) has this representative, and two
degree-zero representatives differ exactly by the denominator
in (C.3bm.2). This proves that formula and the theorem.
\end{proof}

\emph{Exact remaining question.}
For the algebraizable positive-characteristic threefold branch,
the remaining rational top-quotient problem is now the displayed
degree-zero quotient on a projective SNC surface with its explicit
exceptional cycle image. No additional actual-family algebraization
or liftability hypothesis remains on the generators of $J_K(E)$.
The global compatibility of such perfect families across the SNC
intersections is still required. It is not replaced by moving
points independently on components or double curves. The
characteristic-zero branch, higher dimensions without the stated
resolution, and the higher restriction questions remain separate.
